% SOURCE docs/release/2026-10-29/source/docs/causal_banking_release_candidate_2026-10-29.tex CHAR 30046
\chapter{Running and comparing analytical studies}




% SOURCE docs/release/2026-10-29/source/docs/writing-repair/2026-10-29/src/ed71eafe86517078.tex CHAR 0
\section{Commercial orchestration and the recovery of a causal batch}
\label{exec22:workflow-scope}

The chapter compares commercial orchestration, its open clients and delegates, and a distributed execution alternative. Each successive treatment makes the boundary more concrete: a manual describes an operating choice; a client exposes its request and return path; a delegate reveals how retries, state and dependencies are handled. The final Ray comparison asks when distributing the same scientific workload is worth the added recovery and accounting obligations. The evidence supports conditional reuse, with complete scientific outputs and all attempted requests as the common comparison.

A recovered causal batch is useful only if its result answers the declared
question for the declared population and information set. Consider a nightly
offer analysis that fails after writing some cohort effects. Restarting the job
can save computation, but it can also duplicate customers, reuse an outdated
estimate or execute a changed procedure. The relevant product requirement joins
execution recovery to preservation of the scientific request. This motivates
comparing available orchestration before building another scheduler.

The six selected official manuals cover Databricks Lakeflow Jobs on AWS,
Snowflake task graphs, Azure Machine Learning API-v2 pipeline reuse, Dataiku
DSS scenarios and AWS Step Functions. They provide current technical
observations retrieved on 22 October 2026. They do not establish an unchanged
capability at the literature cutoff, 7 October 2026. Databricks' repair and
monitoring pages display 9 and 15 October updates. Azure displays a 30 March
date while its HTML repository metadata records a 19 October update. The
remaining mutable manuals lack an authenticated historical revision. These
observations supplement the cutoff survey; historical coverage remains open.
\citep{exec22_workflow_databricks_repair,exec22_workflow_databricks_monitor,
exec22_workflow_azure_reuse,exec22_workflow_snowflake_graphs,
exec22_workflow_dataiku_scenarios,exec22_workflow_aws_errors}

\subsection{Partial repair must preserve both writes and meaning}

The inspected Databricks repair manual describes rerunning unsuccessful tasks
and their dependents using current job and task settings. Each unsuccessful
task starts again from the beginning. The text explicitly warns that Jobs does
not make task writes idempotent: a partially completed append can duplicate data
when repaired. A shared job cluster is recreated using current settings, and
repair parameters can override earlier values. These mechanisms can support
efficient recovery; they make the application's output-commit and version
contracts material to a causal result. They do not establish that the platform
cannot implement those contracts.~\citep{exec22_workflow_databricks_repair}

Let the finite target population be $\mathcal U$, and let $\tau_i$ be its
stipulated potential-outcome contrast $Y_i(1)-Y_i(0)$. The equal-unit target is
\begin{equation}
 \tau_{\mathcal U}=\frac{1}{|\mathcal U|}\sum_{i\in\mathcal U}\tau_i.
 \label{exec22:workflow-target}
\end{equation}
For an exact constructed example, three units have contrast 6 and one has
contrast $-2$, in hypothetical dollars per unit. Their effect is 4. If the
first two positive-contrast units are written before interruption, and all
four units are appended on repair, the stored-row average becomes
\begin{equation}
 \tau_{\rm append}=\frac{5(6)+1(-2)}{6}=\frac{14}{3},
 \qquad \tau_{\rm append}-\tau_{\mathcal U}=\frac{2}{3}.
 \label{exec22:workflow-duplicates}
\end{equation}
The batch is computationally complete but its implicit positive-group weight
has changed from $3/4$ to $5/6$. At stipulated per-unit cost 1, the intended
net value 3 becomes the misleading value $11/3$. No statistical noise is
needed for this error. These are exact authored values, rather than bank
observations or measured vendor behavior.~\citep{exec22_workflow_reference}

A final output keyed by unit identifier can recover four unique units, provided
the values belong to the same semantic request and the final commit is atomic
or otherwise conflict checked. The history still contains six write attempts.
Deduplicating the result must not erase the partial attempt, its cost or its
failure. Conversely, overwriting a key with a result from a different target
or data vintage is not valid idempotency. Identifier equality alone does not
establish scientific equality. The proposed adapter must distinguish a repeated
delivery of the same result from a conflicting new result.

\subsection{Job status and request status answer different questions}

The current Databricks monitoring manual distinguishes ``Succeeded with
failures'' from ordinary success and explains job outcomes in terms of leaf
tasks. A final task that tolerates an upstream failure can complete without
establishing that every requested customer cohort or sensitivity scenario was
evaluated. A causal batch therefore needs a request-level ledger alongside the
scheduler's result: completed, refused, failed, skipped and reused requests
remain distinguishable.~\citep{exec22_workflow_databricks_monitor}

The same manual describes a 48-hour finished-run graph and a 60-day run-history
window. It also explains that one-time API submissions have no saved job name,
so searching by that name does not find them. A durable scientific record
should bind the actual run identifier and export the relevant lineage before
its operating display expires. Retention and search convenience are useful
platform features; neither determines the validity of the causal assumptions
or the completeness of the requested workload.
\citep{exec22_workflow_databricks_monitor}

For a declared request set $\mathcal Q$, report the requested count together
with each disposition count. Accuracy among returned results and the fraction
of requests receiving an admissible result are different quantities. Retry
success must not silently remove difficult requests from the denominator.
An unsupported intervention is a scientific refusal; a transient storage
failure is an operating failure. Treating them identically would hide the
appropriate remedy and distort the comparison with another route.

\subsection{Version the scientific dependencies behind an execution DAG}

Snowflake's inspected task-graph manual describes dependent SQL execution,
retries, finalizers and overlap controls. Its graph is a scheduling DAG. An
edge that says one job must finish before another does not imply a causal
arrow or identify an intervention. This distinction matters when an agent
translates the same word ``graph'' across workflow and causal APIs.
\citep{exec22_workflow_snowflake_graphs}

The manual says that graph versions include task properties, but not the
definitions of stored procedures called by those tasks. If a called procedure
changes during execution, its new definition can run within the current graph
run. Scientific replay must therefore pin procedure definitions and their
relevant dependencies separately. A stable task-graph version is insufficient
to infer unchanged estimator code.
\citep{exec22_workflow_snowflake_graphs}

The same text says that a suspended child can be treated as though it
succeeded, while a finalizer does not start when the root task is skipped.
Overlap can be allowed where read/write operations tolerate it. These are
documented execution choices whose compatibility with the proposed experiment
must be assessed. A manifest and commit protocol should make incompatible
concurrent writes detectable and should retain skipped scientific requests.
This is a design requirement, not a demonstrated Snowflake adapter or an
absence claim about other product features.
\citep{exec22_workflow_snowflake_graphs}

\subsection{A reusable output must still answer the same target}

Azure Machine Learning's inspected API-v2 description requires a deterministic
component, ForceReRun=false and matching code, environment, inputs, parameters,
output settings and run settings for reuse. The reused component displays the
earlier status, output, logs and metrics. Force-rerun children cannot be reused
by other jobs. These controls offer a useful starting point for lineage and
reproducibility, with the historical-edition limitation stated above.
\citep{exec22_workflow_azure_reuse}

The scientific application must express its dependencies through the declared
artifacts. An unchanged URI need not establish an unchanged table behind it;
an undeclared population revision can change the question even if the code
is identical. The determinism flag is a declaration about the component, not
an independent validation of determinism or causal assumptions. The manual's
code-snapshot and environment comparisons help diagnose execution reuse; the
library must add the estimand and information clock to the meaning of reuse.
\citep{exec22_workflow_azure_reuse}

In the finite example, keeping group effects $(6,-2)$ but changing population
weights from $(3/4,1/4)$ to $(1/4,3/4)$ changes the effect from 4 to 0 and net
value from 3 to $-1$. Retaining the original weights but revising the last
unit's contrast from $-2$ to $-10$ changes the effect to 2 and net value to 1.
The first change is a target change; the second is a data-vintage change.
A cached number may be exactly equal to an earlier output and still answer
the wrong current question.~\citep{exec22_workflow_reference}

\subsection{Retries need a budget that survives redrive}

The current AWS Step Functions manual describes ordered retriers and catchers,
backoff, jitter and bounded retry counts. Catchers apply to Task, Parallel and
Map states, rather than top-level execution failures. Eligible states can have
their retry counts reset on redrive. A library request that requires a total
budget must retain attempts across these boundaries, as well as configure each
state's retry policy. A new execution identifier must not erase earlier time,
resource use or failed scientific requests.
\citep{exec22_workflow_aws_errors}

Error classes require the exact documented meaning. In the inspected current
text, the States.ALL wildcard excludes DataLimitExceeded and Runtime. The
manual explicitly permits a named DataLimitExceeded retry or catch, whereas
Runtime remains non-retryable. It would be inaccurate to infer from the
wildcard exclusion that every explicit payload-limit handler is unavailable.
The historical revision must be recovered before attributing these semantics
to the cutoff edition.~\citep{exec22_workflow_aws_errors}

The manual states that retries are billed as state transitions. Actual
comparison requires more than the number of successful final tasks. For a
common workload, define a proposed accounting identity
\begin{equation}
 C_{\rm workflow}=C_{\rm transitions}+C_{\rm compute}+C_{\rm storage}
                  +C_{\rm transfer}+C_{\rm human}+C_{\rm license}.
 \label{exec22:workflow-cost}
\end{equation}
Each term includes failed and recovery attempts within the declared horizon;
overlapping charges must not be counted twice. Rates, region, service edition,
retention and staff time need authentic evidence before evaluating this
identity numerically. No observed prices, runtime advantage or bank return
are supplied by the exact arithmetic example.

\subsection{What remains unresolved in the commercial comparison}

Dataiku's retrieved page identifies DSS 15 and indexes automation scenarios,
build/check/metrics/code steps, trigger categories and failure/conditional
execution. It is primarily an index. Those headings establish the scope of
linked topics, not detailed transaction, retry or concurrency semantics.
Section~\ref{exec23:dss-scope} develops the technical step, trigger and test controls from separately saved manuals. That resolves the index-only limitation for those mechanisms, while the exact historical server-patch identity remains unresolved.~\citep{exec22_workflow_dataiku_scenarios}

SAS Viya orchestration, Domino and Posit Workbench/Connect receive selected open-client and delegate comparisons in Sections~\ref{exec23:clients-scope} and~\ref{sec:exec24-workflow-contracts}. Those source findings resolve part of this initial candidate inventory; they do not qualify licensed servers, complete capability coverage or supported deployment combinations. Complementary market discovery, product aliases, deployment editions,
licensing and actual prices remain necessary. For the five inspected families,
current manuals support specific mechanisms and application obligations;
none completes historical coverage or native qualification. This comparison
therefore motivates an architecture that reuses suitable orchestration while
adding scientific semantics. Whether each adapter preserves those semantics
remains a proposition to test.

\subsection{A common manifest across the library, CLI and MCP}

Let $m(q)$ serialize the full declared meaning of scientific request $q$:
\begin{equation}
 m(q)=(D,T,I,G,A,E,S,N,O),\qquad k(q)=\operatorname{hash}(m(q)).
 \label{exec22:workflow-manifest}
\end{equation}
Here $D$ identifies input snapshots and availability, $T$ the population and
estimand, $I$ the information clock, $G$ graph/identification assumptions,
$A$ the estimator and covariance settings, $E$ code and environment, $S$ seeds
or randomness coupling, $N$ numerical tolerances, and $O$ the output schema.
The serialization must be canonical and must include the dependencies needed
for that method; the tuple is a proposed contract, not a proof of completeness.
Equal hashes bind specified bytes and declarations. They do not establish
truthful provenance, valid assumptions or an identified causal effect.

The reference program constructs a finite manifest and verifies that changes
to the data, target weights or environment change its key. It also verifies
the append bias and keyed recovery using exact rational arithmetic. This is
a local authored reference, not a vendor adapter or a fault-injected deployment.
The CLI and MCP should transport the same manifest and return lineage, the
target, uncertainty meaning and request disposition. A refused unidentified
effect must not be converted into an ordinary numerical answer by either
entry point.~\citep{exec22_workflow_reference}

\subsection{Demonstrate merit on public fixtures before a costly trial}

A first frozen experiment should combine a public dataset with a separately
generated potential-outcome fixture. The public source exercises realistic
schema, scale and custody; generated outcomes supply known causal truth under
declared assumptions. Public observational agreement alone supplies no such
truth label. Compare uninterrupted and interrupted execution with the same
targets, data, methods, randomization and resource envelope.

Prespecify partial writes, lost responses after commit, changed tables or
procedures, stale cache, skipped children, payload limits and exhausted redrive
budgets. Every requested case remains in the denominator. Deterministic exact
cases require identical results and statuses. Floating or randomized cases
need predeclared tolerances and appropriate coupled randomness. Incompatible
manifests must be refused; commits must be unique and the attempt/cost ledger
complete. A statistical estimator retains a separate calibration protocol.
Computational equivalence does not itself establish coverage or identification.

This experiment can justify whether a native trial is worth its cost. Later
commercial qualification needs the chosen edition, licensed environment,
actual execution and complete prices; later bank acceptance needs authentic
bank questions, truth and owners. The proposed contract supplies a concrete
route for those stages while preserving uncertainty about their outcomes.
It does not establish that any sibling repository or commercial platform has
already adopted the causal library.


% SOURCE docs/release/2026-10-29/source/docs/writing-repair/2026-10-29/src/26b86631b5124aaf.tex CHAR 0
\section{Scenario controls must preserve a fixed scientific workload}
\label{exec23:dss-scope}

The previous orchestration comparison motivates a closer question for a
nightly bank offer analysis: which requested effects survive a partly failed
scenario? Dataiku's detailed official manuals describe useful controls for
handling failures, testing a flow and initiating another run. Their current
technical text supports a more precise comparison than an index of scenario
features. The observations below concern the saved DSS 15 and DSS 14 manuals
acquired after the 7 October 2026 cutoff. They establish documented mechanisms,
with the historical qualification developed below; they do not establish an
executed adapter or an exhaustive account of the product.
\citep{exec23_dss_latest_flow,exec23_dss_latest_triggers,exec23_dss_latest_steps,
exec23_dss_dss14_flow}

\subsection{A tolerated operating error leaves a scientific question}

The step-control manual specifies sequential execution until the first failure
or completion of all steps. Some steps, including Build/train, can downgrade
an error to a warning and let later steps run. An example uses Ignore failure
when attempting to kill another scenario. Individual steps can be disabled
with Never, run with Always regardless of previous failures, or use conditions
on earlier failures, the current outcome or variables. The documented scenario
outcomes are SUCCESS, WARNING, FAILED and ABORTED. These distinctions are useful:
cleanup can still run after a failure, and a later diagnostic need not be lost.
They leave the application responsible for deciding whether its requested
causal effects have actually been returned.
\citep{exec23_dss_latest_flow,exec23_dss_dss14_flow}

Let $\mathcal Q$ be the declared set of cohort requests, $w_q$ their fixed
population weights with $\sum_{q\in\mathcal Q}w_q=1$, $\tau_q$ the cohort's
mean potential-outcome contrast between contact and no contact, and $c_q$
the incremental cost of contacting a member. The declared all-contact policy
has incremental net value, in dollars per member,
\begin{equation}
 V_{\mathcal Q}=\sum_{q\in\mathcal Q}w_q(\tau_q-c_q).
 \label{exec23:dss-target}
\end{equation}
For the following exact authored construction, cohort effects are stipulated
truth, not estimates whose identification or precision has been established.
There are two equal-weight cohorts, with $\tau_A=10$, $\tau_B=-30$ and
$c_A=c_B=1$. No contact has zero incremental value and cost. Thus
\begin{equation}
 V_{\mathcal Q}=\tfrac12(10-1)+\tfrac12(-30-1)=-11.
 \label{exec23:dss-full-net}
\end{equation}
If the effect for cohort B is missing after a failed or skipped step, averaging
the available cohort instead reports $V_{\rm available}=9$. The apparent
improvement of 20 dollars comes from changing the population. It reverses the
financial sign without any sampling noise. A tolerated error that leaves the
same effect missing creates the same arithmetic problem. These are constructed
missing-output illustrations; they are not recorded Dataiku run outcomes.
\citep{exec23_dss_denominator_reference}

The proposed scientific ledger keeps two requested effects and one returned
effect in every such illustration. It retains the failure, skip or tolerated
error as a distinct disposition, the declared weights and the missing result.
It can describe the available cohort's value for that cohort, but cannot
rename it the original population value. Assigning zero to a missing effect
would add a substantive assumption. Even renormalizing costs to the returned
cohort would change the all-contact policy being evaluated.

An explicit range can communicate a limited answer. If a separate substantive
argument justifies $\tau_B\in[-50,50]$, with cohort A's effect and the costs
still stipulated, the original target obeys
\begin{equation}
 V_{\mathcal Q}=4+\tfrac12\tau_B\in[-21,29].
 \label{exec23:dss-bounds}
\end{equation}
This identified set follows from assumed bounds, rather than statistical
confidence. It gives no guarantee of a positive net return. The route to a
sharper answer is to recover the missing effect, defend tighter bounds or
declare a different policy and population explicitly. Reporting only the
successful cohort cannot supply the lost information.
\citep{exec23_dss_denominator_reference}

\subsection{A trigger observes its rule, not a complete causal snapshot}

The trigger manual distinguishes initiation from execution. Several active
triggers are evaluated independently; any true condition can promote a scenario
for execution, which need not start immediately. For filesystem datasets,
change detection uses filenames, lengths and modification times, with an
optional marker file. A dataset-change trigger does not detect SQL data
changes; a SQL-query trigger compares the query's output. The suggested row
count or maximum-column query is a useful proxy, whose adequacy depends on the
scientific request.~\citep{exec23_dss_latest_triggers}

For example, two fixed records can retain both row count and maximum date
while cohort B's stored contrast changes from $-30$ to $30$. In the same
stipulated equal-weight population, net value then changes from $-11$ to $19$.
A count-and-maximum-date proxy cannot distinguish these two tables. This is a
logical counterexample for that chosen query, rather than a claim that SQL
triggers cannot use a stronger query or that the platform loses changes.
The proposed adapter should bind the relevant data snapshot and information
clock, with a change rule appropriate to those dependencies.
\citep{exec23_dss_denominator_reference}

By default, Suppress triggers while running is enabled: a second trigger can
go unevaluated during a long run. If suppression is disabled, an evaluation
finishing while the scenario is still running is ignored. For eligible tracked
events, the manual describes disabling suppression, enabling Delay runs and
disabling Squash delayed runs to retain a run per trigger. It also warns that
serial scenarios accumulate delayed runs if arrival throughput exceeds service
throughput. Squashing limits the delayed queue to at most one run and can
cancel an intermediate run. Time-based, manual and hidden Run another scenario
triggers are excluded from run delaying.
\citep{exec23_dss_latest_triggers}

These controls support different operational objectives. A latest-snapshot
analysis may deliberately combine changes, whereas a predeclared event study
may need every eligible event and its clock. The latter requires a durable
event/request record and reconciliation against the declared workload.
Trigger counts alone cannot certify that denominator, and every raw data
change need not be a scientific event. The application must define the
relationship before choosing suppression or squashing. This is a proposed
contract to qualify, rather than demonstrated event completeness.

\subsection{Existing tests help define a bounded integration experiment}

The scenario-step manual documents checks and data-quality rules: the step
fails when at least one selected check fails, and rule outcomes become
variables for later steps. It also documents project pytest selection and a
flow integration test that swaps in stable reference inputs, rebuilds outputs
and compares them with reference datasets. These mechanisms could support a
prototype comparison using public or generated data. A causal extension would
need reference targets, weights, information clocks, refusals and uncertainty
alongside ordinary output comparisons; the documentation does not show that
such an extension has been executed.
\citep{exec23_dss_latest_steps}

Run another scenario waits for completion and takes the child's outcome.
The manual says the initiating step's user is used regardless of the child
scenario's Run as setting. A proposed adapter must therefore preserve the
observed run identity together with the scientific request, then test the
relevant configuration. This narrow documented behavior does not establish
the whole authorization model. Similarly, a Clear step deletes dataset or
folder contents; persistent scientific evidence must have its own deliberate
retention arrangement.~\citep{exec23_dss_latest_steps}

The economical next experiment is a small fixed manifest containing complete,
failed, skipped and tolerated-without-effect requests. Check the returned
scientific ledger, weights and source identities before measuring retries or
runtime. Pair this with two data versions that share a selected trigger proxy
but differ in a scientific dependency. A successful native experiment would
support only its frozen release, configuration and manifest. Its merit should
be established with public/generated inputs before purchasing a large bank
study. No native execution, licensed entitlement or measured commercial cost
is asserted here.

\subsection{The cutoff requires a patch identity, not just a major version}
\label{exec23:dss-cutoff}

The current release index lists DSS 15.0.0 on 14 August 2026 and DSS 15.0.2 on
24 September 2026, both before the client cutoff. It lists DSS 15.1.0 on
8 October and DSS 15.2.0 on 15 October, both later. DSS 15 therefore cannot be
dismissed as entirely post-cutoff. Equally, its current major-version manual
cannot silently stand for the exact 15.0.2 technical bytes. The DSS 14 release
index lists 14.7.5 on 24 September. Its inspected step-control body matches
the DSS 15 body, which supports a current textual comparison but does not
authenticate the historical patch revision.
\citep{exec23_dss_latest_releases,exec23_dss_dss14_releases,
exec23_dss_latest_flow,exec23_dss_dss14_flow}

Release-index headings for Agent Skills and MCP server do not establish their
detailed behavior. Those claims need the relevant primary technical text and
edition. The detailed scenario review narrows one earlier source omission;
exact cutoff originals for the selected commercial server manuals and complementary market coverage remain to be investigated. The following SAS, Domino and Posit client studies supply selected source contracts and preserve their distinct remaining server and deployment questions. Rights,
prices and native qualification remain separate questions. This comparison
supports reuse of documented orchestration mechanisms as candidates, while
requiring an explicit scientific workload before an integration is accepted.


% SOURCE docs/release/2026-10-29/source/docs/writing-repair/2026-10-29/src/8208ceec7e01cb3c.tex CHAR 0
\section{Commercial execution platforms and open client contracts}
\label{exec23:clients-scope}

A bank's overnight causal analysis may run in an established execution
platform while its scientific meaning remains the responsibility of the
proposed library. A useful commercial comparison therefore follows the
actual boundary: submission, observation of a remote job, interpretation of
its terminal state, recovery and custody of its outputs. The open clients
supplied by SAS, Domino and Posit reveal concrete integration choices. Their
selected source revisions precede 7 October 2026. They authenticate client
and SDK behavior visible in source, not every feature of a licensed server,
a supported deployment combination, or a successful native trial.
\citep{exec23_workflow_sasctl,exec23_workflow_dominopy,
exec23_workflow_connectapi,exec23_workflow_positgo}

\subsection{Completion is one coordinate of a scientific result}
\label{exec23:clients-completion}

Let $q$ denote an immutable scientific request, $a$ a run attempt, $v$ a code
revision, $d$ a data snapshot and $e$ an execution environment. The requested
result has a declared estimand, population, information time and assumptions.
An execution record is a different object: it says which attempt ran which
inputs, whether it terminated, and which artifacts it produced. A proposed
acceptance rule can make these requirements explicit:
\begin{equation}
 A(q,a)=T(a)\,S(a)\,B(q,a,v,d,e)\,E(a)\,C(q,r_a),
 \label{exec23:clients-acceptance}
\end{equation}
where each factor is a zero--one indicator. Here $T$ means terminal, $S$
means an allowed successful execution state, $B$ means a verified binding
between request and execution identities, $E$ means the required result
artifacts $r_a$ exist, and $C$ means the result passes the stated scientific
contract. The last factor may require an explicit nonidentification or
refusal result; a numerical effect need not exist for every valid request.
These are proposed application semantics, not a claim that a vendor implements
this particular rule. A successful process exit supplies evidence for only
part of the conjunction.

This distinction matters even when all arithmetic is exact. Suppose a
funding decision has a stipulated incremental response coefficient $b_v$,
capital exposure of USD 1,000 and cost of USD 20. For two revisions that
encode different supplied assumptions, take $b_0=3/100$ and $b_1=4/100$.
Their net values are
\begin{equation}
 V_0=1000(3/100)-20=10,\qquad
 V_1=1000(4/100)-20=20.
 \label{exec23:clients-revision}
\end{equation}
Both calculations can exit successfully. Substituting the second revision
while describing the run as a replay of the first changes the answer by
USD 10. This is an authored identity example with stipulated coefficients,
not an estimated causal effect or a measured vendor failure. Recording a
code commit alone also leaves the data snapshot, configuration, dependency
versions and randomization state to be specified.
\citep{exec23_workflow_client_reference}

\subsection{Domino: versioned restart and distinct polling predicates}
\label{exec23:clients-domino}

The selected Domino Python client exposes useful controls for this boundary.
Its job submission payload includes a commit identifier, main-repository Git
reference, environment identifier and hardware selection. Its restart method
defaults to requesting the original input commit; the alternative requests
the current version. The stop method separately accepts a flag asking to
commit results. These are meaningful strengths for controlled experimentation,
but the selected client source does not establish complete data and environment
immutability or the licensed server's interpretation of every field.
\citep{exec23_workflow_dominopy}

Two polling interfaces in the same source have different visible predicates.
The legacy blocking run helper resets its consecutive status-read retry count
after a successful read. It waits for an output commit and then requires the
reported status to be \texttt{Succeeded}. The newer blocking job helper
raises on \texttt{executionStatus == Failed} and passes
\texttt{isCompleted} to the polling library as its success predicate. Its
defaults specify a five-second poll interval and a 6,000-second polling
timeout, with selected request exceptions ignored during polling.
\citep{exec23_workflow_dominopy}

These observations justify a precise question rather than an accusation of
a product defect. A supplied hypothetical response whose completion flag is
true and whose execution status is neither Failed nor Succeeded satisfies
the newer local predicate. It would fail a proposed explicit Succeeded
allowlist. This is a logical witness about the inspected expression; it does
not show that the server actually emits such a response. Resolving that
question needs the supported server state schema and its version. An adapter
should preserve both fields, the raw response, job identity and output binding
until the mapping is justified. It must also distinguish failure to observe
a job before the deadline from proof that the remote job has stopped.
\citep{exec23_workflow_client_reference,exec23_workflow_dominopy}

For the proposal, the preferred route is a narrow client wrapper that keeps
Domino's execution and revision controls and adds the scientific request
contract. Reimplementing the entire scheduler has no demonstrated advantage.
The wrapper remains a proposed integration: supported API routes, exception
mapping, actual terminal states, output custody and fault behavior still need
their own inspection and registered public or synthetic trials.

\subsection{Posit Connect: tasks, jobs and schedules have different meanings}
\label{exec23:clients-connect}

The inspected R client declares development version 0.12.1.9000. Its task
object normalizes an incoming identifier when necessary. The
\texttt{poll\_task} implementation waits for the task's
\texttt{finished} field, checks the final numeric code, raises on nonzero
code and attaches an optional result to the task object. The function exposes
a wait interval but no overall deadline parameter. This describes the local
function signature and loop, not the absence of time limits elsewhere in the
server or transport. An integrating worker should impose an explicit
observation budget and retain the remote task identity for reconciliation.
\citep{exec23_workflow_connectapi}

Job records carry additional context: content and bundle identifiers, start
and end times, execution status, exit code and, where available, cluster and
image information. The selected documentation distinguishes Active, Finished
and Finalized on newer servers and identifies missing fields on versions
earlier than 2022.10.0. The implementation also normalizes older timestamp
and identifier representations. A consumer must not treat missing historical
fields as verified empty values or infer a successful scientific result from
the label Finished. A bundle identifier is useful provenance, but its
mapping to the request's code, data and configuration still needs evidence.
\citep{exec23_workflow_connectapi}

Scheduling adds another boundary. The inspected schedule-update method is
marked experimental, converts start and next-run times and calls unversioned
schedule endpoints. Thus the R package's source commit, the server API
behavior and the schedule's information clock are separate versioning
questions. In banking, a job scheduled after a ledger cutoff must not quietly
use records first available after the intended decision time. The proposed
adapter should preserve that time explicitly and reject ambiguous timestamp
conversion. This is an application requirement, not a finding that Posit
violates it.
\citep{exec23_workflow_connectapi}

The R client is therefore a plausible integration route for R-based
estimation and reporting. Its task, job and schedule objects should remain
distinct in the library/CLI/MCP contract. Porting the client wholesale to
another language would add an API-maintenance obligation without yet proving
better scientific fidelity. A worker boundary with typed results is the
more limited candidate to test first.

\subsection{Posit Launcher: reuse the scheduler boundary with explicit state}
\label{exec23:clients-launcher}

The Launcher Go SDK addresses a different layer. Its official guide describes
a plugin connecting Workbench or Connect to a scheduler using length-prefixed
JSON over standard input and output. Bootstrap reconciles current jobs,
heartbeats detect an unresponsive plugin, and multiple Launcher instances
require a common view of job state. The guide gives Go and product-version
prerequisites. These are documented design contracts of the pinned SDK;
they are not a measured availability or recovery guarantee.
\citep{exec23_workflow_positgo}

The testing guide offers a useful structure for proposed qualification. It
separates universal invariants, product request sequences and individual
control scenarios. Examples cover nonempty job identifiers, distinct IDs for
two submissions, retrieval and filtering, terminal-state nonregression,
launch, stop, kill and cancel. Profiles allow scheduler-specific terminal
states and exit codes; unsupported suspend/resume can be skipped. An
integration must interpret the selected profile and skipped cases rather
than collapse the entire suite to one portable success claim. These tests
are documented upstream facilities; no passing execution is asserted here.
\citep{exec23_workflow_positgo}

For a bank with an accepted scheduler, adopting this bridge boundary is more
plausible than reproducing an execution platform inside a causal library.
The library would still own causal request identity and result validation.
The platform would own job dispatch and resource management. A prototype
comparison should test stale state after restart, duplicate submission,
status-stream interruption and missing output, including the cost of state
reconciliation. A failed bootstrap or missing shared state must remain
visible as an operating failure, not be translated into a zero causal effect.

\subsection{SAS: preserve experimental scope and explicit timeout outcomes}
\label{exec23:clients-sas}

The sasctl project exposes Python module and command-line access to SAS
Viya. Its selected Workflow class states that its partial implementation is
experimental and can change without notice. The comment refers to the
documentation situation in September 2022; it does not establish whether
public workflow documentation exists at the proposal cutoff. Submission
selects an enabled definition, maps supplied values to declared prompts
and includes their declared types and versions. Naive datetime values are
converted using local timezone when possible, with a UTC fallback if they
remain naive. A banking adapter should require an explicit timezone before
submission rather than let the worker host silently define the information
clock.
\citep{exec23_workflow_sasctl}

A separate score-execution method exposes another concrete control issue.
It returns the strings completed or failed on those reported states. Its
timeout test is inside the polling loop, while the loop guard itself limits
elapsed time; there is no explicit return after the loop. A supplied zero
timeout therefore skips the body and reaches Python's implicit None return.
For a positive budget, the timeout branch can be reached if time crosses
the threshold during an iteration, but a later failed loop guard can also
end the loop without that explicit timeout result. This is source-level
control-flow reasoning, not a timed experiment against SAS.
\citep{exec23_workflow_sasctl,exec23_workflow_client_reference}

The proposed wrapper should validate a positive observation budget, normalize
all completion paths into typed states, retain execution identity and logs,
and require a separately verified result binding. An unrecognized or absent
polling return is an unresolved observation, not permission to use stale
scores. Reusing the open client selectively is a candidate; depending on its
experimental workflow service without versioned qualification is not yet a
supported product commitment.

\subsection{Retries, rights and a costed prototype choice}
\label{exec23:clients-choice}

Recovery can change value when a retry repeats a charge or an externally
visible action. In an exact illustrative servicing example, suppose one
useful contact produces a stipulated USD 12 gain and every attempted
delivery incurs USD 7. If there is at most one useful gain, net value after
$n$ charged attempts is
\begin{equation}
 V(n)=12-7n,\qquad V(1)=5,\qquad V(2)=-2.
 \label{exec23:clients-retry}
\end{equation}
The sign changes solely because of the assumed duplicate cost. A repeated
analysis job need not incur this contact cost; the example concerns a
workflow that crosses into a paid or externally visible action. It motivates
separate request and attempt identities, reconciliation before resubmission
and an idempotency key at the side-effect boundary. It establishes no vendor
failure rate, bank return or economic superiority.
\citep{exec23_workflow_client_reference}

The selected connectapi license is MIT; the Launcher root license provides
an MIT-style grant unless supplied under a commercial agreement. The Domino
and sasctl client roots include Apache 2.0 texts. These source licenses are
relevant to proposed client reuse. They do not supply an entitlement to a
commercial server, a supported edition, a complete dependency-rights review
or a price. Consequently the commercial comparison should price an actual
deployment and support arrangement separately from the open wrapper.
\citep{exec23_workflow_connectapi,exec23_workflow_positgo,
exec23_workflow_dominopy,exec23_workflow_sasctl}

The practical proposal is to compare narrow wrappers around established
execution services using the same public or synthetic scientific workload.
Record installation and maintenance effort, queue and compute time, output
validation, failed attempts, recovery and any billed actions. Preserve every
requested case in the denominator. Qualify client and server versions,
explicit failure states, output custody and clock semantics before drawing
an economic conclusion. Historical Dataiku technical-version gaps and wider
market coverage remain separate survey obligations; these four inspected
profiles neither establish market completeness nor erase those limits.


% SOURCE docs/release/2026-10-29/source/docs/writing-repair/2026-10-29/src/dde61b4e3ec047b3.tex CHAR 0
\section{Following commercial workflow requests into their delegates}
\label{sec:exec23-workflow-delegates}

An overnight funding estimate is usable only if its execution record still
refers to the requested population, information date and model. The boundary
between a causal library and an execution platform therefore includes more
than a submission endpoint. Transport retries, authentication refresh,
pagination, scheduler state and returned artifacts can each change what a
consumer believes happened. Following the pinned SAS, Domino and Posit
clients into their delegates makes the comparison more precise. These are
selected primary files from source revisions before 7 October 2026, rather
than observations of paid servers or a complete server compatibility matrix.
The preceding client comparison remains applicable in its stated scope.
\citep{exec23_workflow_delegates_sasctl,
exec23_workflow_delegates_dominopy,exec23_workflow_delegates_positgo}

\subsection{SAS: separate refresh counts, transport time and sampling membership}

The SAS service boundary is concrete. \texttt{Service.request} prefixes its
service root and forwards the request and keyword arguments to
\texttt{core.request}; the latter calls the selected session and raises an
HTTP error for response codes 400 through 599 before converting the response.
The session exposes \texttt{timeout=None} and forwards the value to its
underlying request implementation. When a 401 challenge specifically reports
an expired access token, successful refresh can cause the request to be
issued again. This supports a wrapper using the existing client while making
request identity, authentication events and transport policy explicit. It
does not demonstrate that every failed submission is safe to repeat or that
a remote side effect occurred twice.
\citep{exec23_workflow_delegates_sasctl}

The generic \texttt{Service.\_monitor\_job} helper illustrates a second
boundary. It treats both completed and failed as terminal, sleeps half a
second before each refresh, and limits the number of refreshes. On a
terminal response it returns the job, including a failed job; otherwise it
raises \texttt{JobTimeoutError}. That is a count-limited observation
procedure. It does not itself establish a wall-clock deadline for a request
that stalls. The inspected workflow-submission method is not shown to call
this generic helper, so this finding must not be projected onto every
workflow execution path. The workflow unit tests mock definition lookup and
POST: they establish intended URL and prompt formatting, including conversion
of datetime prompts to strings, rather than actual server completion.
\citep{exec23_workflow_delegates_sasctl}

For a supplied timing example, let $m$ be the number of nonterminal
refreshes, $h$ their sleep interval and $r_i$ the duration of refresh $i$.
Ignoring fixed overhead, the observation time is
\begin{equation}
 t_m=mh+\sum_{i=1}^{m}r_i.
 \label{eq:exec23-delegate-time}
\end{equation}
With three refreshes, $h=1/2$ second and every $r_i=2$ seconds, the sleep
budget is $3/2$ seconds while the actual supplied duration is $15/2$
seconds. An upper bound on $m$ alone cannot bound the sum if request
durations are unrestricted. These are exact authored timings, not measured
SAS latency. A proposed worker should combine a monotonic observation
deadline with bounded transport requests and retain the remote job identity
when the deadline expires. Expiration establishes that observation stopped;
remote termination requires separate reconciliation.
\citep{exec23_workflow_delegates_reference}

Pagination also matters to scientific meaning. The selected
\texttt{PageIterator} extracts start and limit from a next link and schedules
offset requests through a thread pool. It stops after an empty result. The
item iterator explicitly calls its length an approximate count and notes
that permissions can affect accessible membership. This offers convenient
client iteration without proving a frozen population, stable ordering or
snapshot consistency. It would be incorrect to infer that the server lacks
such facilities merely because the selected iterator does not establish
them; the server contract must supply the missing evidence.
\citep{exec23_workflow_delegates_sasctl}

Consider an authored frame with identifiers $(1,2,3,4)$, page size two and
binary outcomes $(1,0,1,0)$. The first page returns $(1,2)$. If identifier 1
is removed before a second request with offset two, the current frame is
$(2,3,4)$ and the second page returns $(4)$. The consumer observes
$(1,2,4)$ and estimates $1/3$, whereas the initial-frame mean is $1/2$.
This is a changing-offset counterexample, not a reproduced SAS defect.
Deduplicating returned IDs would not recover the missing initial identifier
3. A causal audit that samples workflow records must therefore bind the
membership and permission law to an immutable snapshot, a reconciled ID
manifest or a justified server cursor contract. More parallel requests do
not by themselves identify that law.
\citep{exec23_workflow_delegates_reference}

The preferred SAS route remains a narrow Python/CLI worker over declared
service versions. Reusing the request and prompt machinery avoids creating
another partial REST client. Wholesale porting would still require the
server schemas, time conventions, authentication and supported edition
contracts that are currently unresolved. A future registered prototype
should distinguish an accepted submission, an authentication refresh, a
failed job, an observation deadline and a confirmed cancellation. No native
trial or server entitlement follows from the open client source.

\subsection{Domino: preserve method policy and exception identity}

The Domino route delegate distinguishes legacy run endpoints from the v4
job endpoints for start, stop, status, restart and execution details. Its
HTTP manager mounts an HTTPS retry adapter with a configured total whose
default is four, backoff factor one and status list 408, 502, 503 and 504.
It does not explicitly override the retry library's allowed-method policy.
Consequently, the presence of this adapter is insufficient evidence that a
job-submission POST is automatically retried on those statuses. That claim
needs the exact HTTP dependencies and their policy. The wrapper forwards
request keywords without setting an overall deadline in these methods;
the blocking helper's polling timeout is a separate layer.
\citep{exec23_workflow_delegates_dominopy}

Exception classes also affect what the polling loop does. The manager
converts HTTP 409 into \texttt{ReloginRequiredException}; other HTTP errors
are re-raised. The selected exception inherits \texttt{DominoException},
whose base is \texttt{Exception}. It is not the
\texttt{requests.exceptions.RequestException} listed in the blocking
helper's default ignored exceptions. Thus the statement that this helper
ignores intermediate request exceptions must retain its declared class
boundary. It cannot be expanded into a claim that authentication-expiry
exceptions, every HTTP conflict or every failure is silently retried. The
409 mapping is a client interpretation; an authoritative server error
contract is still needed before assigning the same meaning to every
supported deployment.
\citep{exec23_workflow_delegates_dominopy,exec23_workflow_dominopy}

The job tests clarify intent without closing that deployment question.
Unit fixtures provide Queued and Succeeded responses, an
\texttt{isCompleted} flag and a mock server version 9.9.9. Other unit tests
specify a polling timeout when \texttt{RequestException} is ignored, and
immediate propagation when the ignored-exception tuple is empty. Live tests
are conditional on a reachable Domino deployment. Reading these sources is
neither running the tests nor authenticating 9.9.9 as a real supported
server. The earlier question about a complete terminal-state schema remains
open: a successful mock does not enumerate every server terminal state.
\citep{exec23_workflow_delegates_dominopy}

For the proposal, preserve the client's useful restart and revision controls
and add an explicit result contract at the worker boundary. Keep transport
retry, reauthentication, remote restart and a new scientific attempt as
different events. Before any resubmission with external effects, reconcile
the remote attempt and use an application identity whose idempotency
semantics are actually supported. This is a proposed protocol, not a claim
that Domino duplicates jobs or lacks an idempotency facility. Porting the
whole client would reproduce the unresolved schema and dependency questions
without yet demonstrating better causal fidelity.

\subsection{Posit: terminal state, exit code and conformance coverage}

The Launcher SDK gives a particularly useful typed distinction. Its
\texttt{Failed} status means the job could not be launched; its
\texttt{Finished} status includes a process that launched and exited with
a nonzero exit code. Canceled, Failed, Finished and Killed are terminal.
The job object carries status, optional reason/code and a separate optional
process exit code. A causal CLI or MCP consumer should preserve these fields
instead of translating every terminal or Finished record into scientific
success. A successful refusal artifact may be an acceptable scientific
answer, but that acceptance must be explicitly bound to its request and
contract.
\citep{exec23_workflow_delegates_positgo}

To expose the accounting consequence, suppose four supplied attempts yield
Finished with exit zero and a valid bound artifact; Finished with exit seven;
launch failure; and cancellation. All four are terminal. Only the first
satisfies a stipulated acceptance predicate requiring the allowed state,
allowed exit and artifact binding. At a supplied cost of USD 3 per attempt,
USD 12 has been spent. Reporting USD 3 per terminal event and USD 12 per
accepted artifact answers different questions. Neither number proves causal
validity, which remains another contract, and neither is a Posit price or a
bank ROI estimate. This exact example shows why a proposal comparison must
use the same accepted-output definition and failed-request denominator for
every execution alternative.
\citep{exec23_workflow_delegates_reference}

The scheduler example offers a reusable separation between Launcher protocol
handling and a \texttt{SchedulerClient} backend that submits, observes,
cancels, kills and retrieves output. Its illustrative state mapper sends
unrecognized native values to Pending. That default is an example, not
evidence that all scheduler states are known or that a supplied plugin is
production ready. A bank worker should retain the native value and use an
explicit unresolved mapping until the chosen backend profile justifies it.
Reusing the SDK boundary is a narrower candidate than reimplementing the
execution service, but the backend and supported product versions must be
selected before qualification.
\citep{exec23_workflow_delegates_positgo}

The source conformance facilities strengthen the proposed qualification
design when their coverage is kept explicit. A profile requires job
factories and expected stop/kill states and exit codes. Network, output and
suspend switches determine which workflow assertions run. The
\texttt{NetworkAvailable} comment says default true, while the inspected
\texttt{applyDefaults} function sets only durations. An omitted Go boolean
remains false. A prototype must therefore record the actual profile and set
the required switch explicitly, instead of assuming that a documented
default caused a network check to execute. This is a discrepancy between
selected source comments and defaults, not an observed deployment failure.
\citep{exec23_workflow_delegates_positgo}

Two further qualifications prevent overstatement. The terminal-regression
invariant compares a terminal result with one revisit 200 milliseconds
later. Agreement at those times is useful evidence for that run, not a proof
of permanent stability: an authored history can agree at times zero and
one fifth second and change at one second. The cancellation scenario skips
when the API returns an error, explaining that the job may already have
started. That branch does not first demonstrate the proposed start race.
A registered adaptation should inspect state and error identity before
classifying a race, retain every skipped or unsupported case, and avoid
counting a skip as successful cancellation. No upstream conformance suite
has been executed here.
\citep{exec23_workflow_delegates_positgo,
exec23_workflow_delegates_reference}

\subsection{A bounded integration choice and the remaining comparison}

These delegates support a coherent prototype choice. Preserve vendor job
dispatch and resource management behind typed workers, with a common causal
request and result contract above them. Record request, attempt, code,
dataset, environment, information date and native state separately. Bound
observation and transport time; reconcile a remote attempt after a timeout;
freeze audit membership; and validate the expected artifact and its
scientific meaning independently of process completion. This is a proposed
integration architecture. The exact authored timing, sampling and cost
witnesses explain why its contracts are needed; they do not qualify the
implementation.

The remaining market comparison has specific inputs. The next treatment adds SAS's historical REST examples, a declared-range HTTP dependency profile for Domino and Posit's concrete Local backend. Those resolve specific source questions. SAS still needs a complete supported server contract, Domino its terminal and error schema, and Posit a supported product/backend combination with an explicit conformance profile. Immutable Dataiku 15.0.2 cutoff technical bytes remain a
separate historical-source omission. Client licensing, commercial server
rights, dependency notices and actual prices must be assessed separately.
Public or synthetic fault trials can justify a later costly prototype study
once its environment is fixed. Until then, these are source-based comparison
and integration choices; native capability, whole-market coverage and bank
acceptance remain open.


% SOURCE docs/release/2026-10-29/source/docs/writing-repair/2026-10-29/src/ea6f7e39725914f9.tex CHAR 0
\section{From workflow submission to a scientifically acceptable artifact}
\label{sec:exec24-workflow-contracts}

A restarted causal batch needs an account of what was requested, what
actually ran and which result answers the banking question. A collections
team might request a monthly policy comparison, receive a trigger
acknowledgement, and later find either a completed simulation or no run at
all. A funding analyst might receive a Finished job whose process exited
unsuccessfully. Neither an acknowledgement nor a terminal state establishes
that the requested intervention contrast, population and information date
survived execution. The following historical source profiles supplement the
earlier client comparisons with immutable Dataiku client text, SAS REST
examples, a declared-range HTTP dependency profile and a concrete Posit
Local backend. They establish source contracts, not observed performance of
licensed installations.~\citep{exec24_workflow_contracts_dataiku,
exec24_workflow_contracts_sas,exec24_workflow_contracts_domino,
exec24_workflow_contracts_posit}

\subsection{Dataiku: the trigger denominator precedes the run denominator}

The authenticated Dataiku client identifies version 14.7.3. Its scenario
interface explicitly distinguishes a trigger firing from a scenario run.
A manual firing may be cancelled because a scenario is already running or
another trigger fires. The firing's run identifier is the trigger-run
identifier; it is not automatically the eventual scenario-run identifier.
The proposed adapter should retain both identities and their relation,
including the legitimate outcome that no scenario run is created. This
historical client cannot authenticate the separately requested DSS 15.0.2
server behavior or a complete supported client/server matrix.
\citep{exec24_workflow_contracts_dataiku}

The distinction matters even in convenience methods. The inspected
trigger-wait method with \texttt{no\_fail=True} returns \texttt{None}
when a trigger is cancelled. The surrounding \texttt{run\_and\_wait}
method passes that option into the trigger wait, then constructs a waiter
without checking whether a run exists. The waiter's next loop dereferences
\texttt{scenario\_run.run}. Under the documented cancelled-trigger
condition, this composition therefore has an unguarded \texttt{None} path.
This is a conditional static finding; no native execution, deployment
prevalence or observed customer failure is claimed. A wrapper should handle
a cancelled-before-start event explicitly instead of promising that the
permissive flag suppresses every exception. It can do so without porting
the whole scenario engine.~\citep{exec24_workflow_contracts_dataiku}

Once a run exists, the client treats the presence of its result as the end
of running. It exposes SUCCESS, WARNING, FAILED and ABORTED outcomes.
The strict completion path accepts only SUCCESS; WARNING is not silently
equivalent to it. A permissive return still needs an explicit consumer
decision about outcome and artifact validity. Backoff begins at 200
milliseconds, grows by a factor of 1.1 and caps each sleep at 15 seconds.
This per-sleep ceiling does not bound the number or duration of requests.
The inspected DSS HTTP delegate does not pass a timeout argument. These
facts motivate a separate observation deadline and transport policy; they
do not establish that every customized client session behaves identically.
\citep{exec24_workflow_contracts_dataiku}

For an authored example, suppose four requested triggers produce one
cancellation before a run, one successful accepted result, one warning
result outside the acceptance contract, and one failed run. There are four
requests, three actual runs and one acceptable artifact. If $A$ is the
number accepted, $N$ the number requested and $R$ the number started,
\begin{equation}
  \widehat p_{\mathrm{request}}=A/N=1/4,
  \qquad \widehat p_{\mathrm{started}}=A/R=1/3.
  \label{eq:exec24-workflow-denominator}
\end{equation}
Both rates are legitimate when labelled, but only the first retains the
cancelled request. With stipulated costs of USD 2, 4, 3 and 3, total cost
is USD 12: USD 4 per started run and USD 12 per accepted artifact. These
invented amounts demonstrate a denominator effect, not vendor prices or
bank economics. Three stipulated backoff sleeps total $331/500$ seconds.
With three two-second requests, observation takes $3331/500$ seconds before
other overhead. A sleep schedule is not a whole-workflow time guarantee.
\citep{exec24_workflow_contracts_reference}

The same distinction protects statistical evidence. Public-data evaluations
that retain only successful runs describe performance conditional on
execution and acceptance. They cannot represent all requested cases without
accounting for exclusions. The proposed CLI/MCP result should expose request,
trigger and run identities, whether execution began, native outcome,
exception and acceptance decision. A caller can then require strict success
or inspect a warning while retaining the original denominator.

\subsection{SAS: an execution example is not a supported server schema}

The pinned SAS repository describes its contents as preview REST examples.
Its Job Execution text describes asynchronous execution of a job definition,
submission through \texttt{/jobExecution/jobs}, and polling the created
job's state link. It distinguishes completed from failed and directs the
consumer to an error member on failure. Artifacts occupy a provider-specific
results map; the log location is separate. Without a supplied expiration,
the example assigns cleanup to the client. A workflow-origin submission
can identify its submitting application. These details support an adapter
that follows resource links and preserves provider identity, execution
state, artifacts and lifecycle actions separately.
\citep{exec24_workflow_contracts_sas}

The page identifies document version 5 and an update date of 26 November
2019, not a SAS server release. A cutoff repository revision authenticates
the bytes without making the examples a complete supported API schema.
The separate Visual Investigator workflow text concerns BPMN case management
for investigation and threat analysis. It must not be substituted for the
analytics workflow service used by an earlier Python client. Supported
editions, complete state/error enumerations, prompt schemas and timezone
behavior still need normative technical originals.
\citep{exec24_workflow_contracts_sas}

For the proposal, a narrow worker can bind the scientific target and data
revision to the definition, submission, provider and returned artifact
before remote cleanup. Reimplementing the REST stack would not remove
the missing server evidence. Public examples justify prototype interfaces
and test cases; they cannot supply licensed-deployment conformance results.

\subsection{Domino: method and failure class govern retry behavior}

The earlier Domino HTTP manager supplies a retry policy without overriding
its allowed-method set. Package setup permits Requests at least 2.4.2 and
urllib3 at least 1.26.19 but below 3; the requirements file narrows Requests
to the 2.31 line. A selected historical profile of Requests 2.31.0 and
urllib3 1.26.20 fits these inspected ranges, including Requests' own range.
It is a source-review profile, not a solved environment or deployment lock.
A separately acquired cutoff development revision cannot replace this
declared-range distinction.~\citep{exec24_workflow_contracts_domino}

Requests retains the supplied Retry object and passes it to the connection
pool. In the selected urllib3 profile, default allowed methods are HEAD,
GET, PUT, DELETE, OPTIONS and TRACE. POST and PATCH are absent. The response
retry predicate checks the method before its status list; the read-error
path also applies a method gate. Thus a retryable status list alone does
not imply that a submission POST is retried on those statuses. The
connection-error branch differs: it does not apply that gate. The source
classifies connection failures as cases in which the server is thought not
to have received the request, whereas read errors can follow processing.
A blanket claim that POST can never be retried would also be too broad.
\citep{exec24_workflow_contracts_domino}

A transport retry, authentication refresh, restart of a known remote job
and new scientific attempt must therefore remain distinct. After an
ambiguous submission, the proposed worker retains the unresolved remote
identity and reconciles it under an actual server idempotency contract.
HTTP defaults and response codes do not prove exactly-once execution.
The pinned public \texttt{api-docs} repository does not close this gap:
its original README and index redirect to a mutable latest guide rather
than storing a versioned \texttt{v4/jobs} schema. Complete terminal/error
semantics and the meaning of every conflict response remain unresolved.
\citep{exec24_workflow_contracts_domino}

The open client remains a plausible integration boundary because it already
exposes identities and job actions. A wholesale port would create a second
implementation with the same unresolved server assumptions. A prototype
should first lock full dependencies and register response-status,
connection, read, credential and reconciliation cases. No native trial or
claim that an installation duplicates or loses requests is made here.

\subsection{Posit: preserve a concrete backend's unsupported actions}

The C++ SDK's sample Local backend supplies a concrete execution contract.
It launches an asynchronous child, stores its process identifier and
publishes Pending. On exit it records the process exit code, then changes
every non-killed job to Finished without requiring zero. A quickly exiting
Pending job can first receive a Running notification. This backend explicitly
does not support cancellation. Kill, stop, suspend and resume use separate
signals and state paths. A proposed adapter must report unsupported Cancel
explicitly instead of turning it into a successful or merely skipped test.
\citep{exec24_workflow_contracts_posit}

Two stipulated records can both say Finished while one has exit zero and
a valid artifact and the other has exit seven and no acceptable artifact.
Acceptance needs the appropriate native state, an allowed exit code where
the contract uses one, a request-bound artifact and scientific validation.
A valid refusal may itself be an acceptable scientific answer if the contract
declares that possibility. The example describes invented records, not
observed bank jobs.~\citep{exec24_workflow_contracts_reference}

Historical documentation limits remain material. The prerequisite guide
mentions Workbench 1.4, later says 1.4 or greater, and requires a launcher-enabled
license and test scheduler. That statement is not a current support matrix.
The build specifies C++11, default Boost 1.83.0 and OpenSSL 1.1.1, which need
reconciliation with the older prerequisite table before building an environment.
The Local three-second setting concerns inter-node connections, not a global
causal-job deadline. The SDK notice distinguishes public permissive terms
from software received under a commercial agreement; it does not license
Workbench itself.~\citep{exec24_workflow_contracts_posit}

The smoke-test guide explicitly says its utility does not replace integration
testing with Launcher and the consuming product. It varies privilege and
configuration and requires manual verification of printed responses. Its
informal description of Finished as without failure must not erase the
executable exit-code distinction. A future profile should name exact product
and Launcher builds, SDK revision, backend, privilege mode, network/output
switches and allowed exit codes. Unsupported Cancel, races, missing artifacts
and failures each retain their own accounting. This is a testable prototype
plan, not completed conformance.~\citep{exec24_workflow_contracts_posit}

These comparisons favor small typed workers around existing execution
contracts. The causal library owns the intervention target, assumptions,
observation clock, artifact validation and failure accounting; platforms
retain their native execution vocabulary. Public evidence now supports a
more precise trigger/run distinction, retry classification and exit-code
contract. Historical server manuals, complete schemas, supported product
profiles, rights and measured costs still need their stated evidence. These
limits guide the next prototype without claiming complete commercial-workflow
coverage.


% SOURCE docs/release/2026-10-29/source/docs/writing-repair/2026-10-29/src/1a7ce759f93f2ca5.tex CHAR 0
\section{Dataiku as a workflow reference with explicit scientific acceptance}
\label{sec:exec25-dataiku-reuse}
Dataiku supplies a concrete orchestration reference for the proposed causal
library. The selected evidence concerns the Python API client at revision
\texttt{de83b58cc2d5e1b539fd}\allowbreak\texttt{44797778d56e20f1d367},
committed 3 August 2026, whose package declaration is version 14.7.3. It also
includes retained official DSS 14 and DSS 15 documentation. The immutable
client predates the 7 October 2026 cutoff; the subsequently acquired mutable
documentation is qualified by branch and acquisition date, and does not prove
an exact cutoff server patch or current commercial offer. The reference choice
separates the open client, its dependencies and the commercial server profile.
A client license notice is not a server entitlement, compatibility guarantee
or causal engine qualification.~\citep{exec24_workflow_contracts_dataiku,exec23_dss_dss14_flow,exec23_dss_dss14_releases,exec23_dss_latest_flow,exec23_dss_latest_releases}

This treatment consolidates the preceding Dataiku source findings into an adoption choice; it is not an additional native trial. A useful banking application is a recurring branch-policy analysis. A scenario
can coordinate preparation, fitting, diagnostics and delivery, but a completed
scenario cannot alone establish that the treatment definition, observation
clock, adjustment set and uncertainty calculation are appropriate. Suppose a
monthly analysis is replayed after a data correction. The proposed library
must bind the question version, eligible branch-periods, treatment adoption
record, data and schema digests, method configuration and output identity.
An output from the previous data version is not an accepted answer to the new
request merely because the scenario has a successful execution status. This
scientific acceptance layer is a proposed responsibility of the causal library,
not a claim that the orchestration product lacks a particular undocumented
feature.~\citep{exec25_dataiku_choice}

\subsection{Trigger, run and result are separate objects}
The selected client distinguishes firing a trigger from obtaining a scenario
run. Its \texttt{run} method returns a trigger-fire object; the waiting path
polls for the run or cancellation. Cancellation can return \texttt{None} when
\texttt{no\_fail=True}. In the inspected \texttt{run\_and\_wait} composition,
that result is passed to a waiter without a local \texttt{None} guard, and the
waiter dereferences the scenario-run object. This identifies a conditional
static path to qualify and handle in an adapter; no native cancellation trial
or observed server failure is claimed. The proposed adapter should keep
\emph{trigger cancelled before run creation} distinct from \emph{run aborted}
and from \emph{request completed with a scientific refusal}.~\citep{exec24_workflow_contracts_dataiku}

Completion and outcome also differ. One inspected path regards the presence
of a result as completion, while its strict outcome check accepts
\texttt{SUCCESS} and raises for other outcomes. The documented outcome set
includes \texttt{WARNING}, \texttt{FAILED} and \texttt{ABORTED}; the selected
waiter can return non-success when \texttt{no\_fail} is enabled. Official DSS 14
step-control documentation permits selected failures to be downgraded to
warnings and provides conditional steps, including steps that run regardless
of earlier failure. A proposed causal adapter must retain both the original
outcome and the configured exception policy. A warning permitted by workflow
configuration may still be scientifically disqualifying, while a deliberately
completed refusal can be a valid service response. Neither should be silently
recoded to a successful effect estimate.~\citep{exec24_workflow_contracts_dataiku,exec23_dss_dss14_flow}

The orchestration boundary must therefore expose a logical request identifier,
trigger identity, run identity and attempt identity separately. Consider an
authored example with requests A, B and C: A has two attempts and two created
runs, B has one attempt cancelled before a run, and C has one attempt and one
run. There are three logical requests, four attempts and three runs. The same
run count as request count is accidental; reporting only that count hides both
A's duplicate execution and B's cancellation. No idempotence guarantee is
inferred from the selected client. The proposed replay contract reconciles
existing run identities and accepted artifacts before retrying, and records a
new execution attempt without creating a second scientific request.~\citep{exec25_dataiku_choice}

\subsection{A bounded observer does not prove a stopped computation}
The selected backoff helper starts at 200 milliseconds, multiplies by 1.1 and
caps each sleep at 15 seconds. A cap on an individual polling delay does not
bound the sum of delays. Once the cap is reached, $n$ further waits consume
$15n$ seconds before any HTTP or server time is added. The inspected HTTP call
does not supply a timeout argument, so this source path does not establish a
whole-operation deadline either. These observations are restricted to the
selected source; they are not an assertion about all server or deployment
controls.~\citep{exec24_workflow_contracts_dataiku,exec25_dataiku_choice}

The proposed contract distinguishes a caller's observation deadline from a
confirmed server cancellation. In a stipulated timeline, a trigger fires at
second 0, a run is created at second 3, observation stops at second 10 and the
run completes at second 12. At second 10 the honest state is completion
unknown. An automatic retry could overlap the original run unless the adapter
reconciles its identity and effects. Conversely, a server acknowledgement of
a cancellation request is not accepted as proof that every side effect has
been rolled back. The candidate integration needs bounded client calls,
monotonic total budgets, authenticated terminal-state reconciliation and an
explicit policy for late results. The library must be able to decline a stale
artifact while preserving the costs incurred by the original attempt.~\citep{exec25_dataiku_choice}

\subsection{Conditional integration and the public qualification route}
The present choice is reference use of the selected client and documented
scenario semantics. Conditional integration would wrap a supported and
entitled server/client combination, retain the distinct states above, and
connect completed artifacts to the library's causal question contract. Direct
adoption as a causal engine is not established by orchestration evidence.
Porting selected client logic requires a component and dependency rights
review and reproduces the responsibility for protocol compatibility; porting
proprietary server internals is not authorized by the client notice. Rejection
is specific to unsupported semantic shortcuts, such as equating a successful
scenario with identified causal effect, or treating a local observation timeout
as confirmed remote termination. Other product functionality is outside this
selected-source assessment.~\citep{exec24_workflow_contracts_dataiku,exec25_dataiku_choice}

Before spending on a bank deployment, a public qualification study should
freeze a permitted server profile, compatible client, project and scenario
versions, generated data law, effect query, artifact digests, acceptance rules
and complete cost categories. Generated branch-period fixtures can carry known
causal truth; permitted public datasets can exercise data preparation and
workflow without supplying unobserved assignment truth. The exact scientific
oracle must be supplied by the selected causal method, separately from the
orchestration checks. Fault fixtures include a trigger cancelled before run
creation, a warning-producing step, late completion after an observer deadline,
a stale output digest, duplicate delivery, a failed diagnostic, a scientifically
valid refusal and a changed question version. A fixed direct-library workflow,
the permitted upstream scenario and the proposed adapter should receive the
same scientific task and information. No trial in this protocol has been run
as part of the reference decision.~\citep{exec25_dataiku_choice}

Scientific owners would approve the query, assumptions and acceptance rule;
workflow owners would maintain scenario-to-artifact identity; platform and
procurement owners would establish server rights, compatibility and support;
evaluators would control the fault and causal oracles; maintainers would
review version drift. These remain proposed responsibilities. An illustrative
USD 1,200 preparation cost, USD 1,800 for all attempts, USD 2,400 for scientific
review and USD 600 maintenance allocation total USD 6,000. Across 120 attempted
requests, 90 completions and 60 admissible outputs, the corresponding costs
are USD 50, USD $200/3$ and USD 100. These authored amounts establish accounting
semantics, not a quote or return on investment. The commercial decision requires
actual entitlements, measured complete costs and accepted causal outputs from
the selected supported deployment.~\citep{exec25_dataiku_choice}


% SOURCE docs/release/2026-10-29/source/docs/writing-repair/2026-10-29/src/7e83a6581138f079.tex CHAR 0
\section{Ray as a conditional execution substrate: preserving the scientific request}
\label{sec:exec24-ray}
A bank adopting distributed computation must preserve the scientific request when a worker is rescheduled, a result is reconstructed or a driver disappears. Ray is a candidate execution substrate for repeated nuisance fits, simulation replications and score computation. It does not define the causal estimand, sampling unit, admissible information set or covariance law. Its usefulness therefore depends on an application contract that survives changes in task placement and retry history. The appropriate proposal decision is to retain Ray as a source-grounded reference and conditional prototype option until profiling and conformance evidence justify distributed adoption.~\citep{exec24_ray_sources,exec24_ray_contract}

\subsection{Artifact identity and the scope of the proposed use}
The inspected repository is pinned to a commit dated 7 October 2026 at 15:33:11 UTC, before the end of that day in Asia/Shanghai. Its version file declares 3.0.0.dev0 and a commit placeholder replaced when wheels are built. A source revision is consequently not an installed release, and the development version must not be presented as evidence that a selected wheel contains these exact APIs. The setup file declares Python at least 3.10 and lists several supported minor versions, base packages and optional groups. Those declarations do not resolve the complete build, native ABI or transitive dependency inventory. The direct license contains Apache 2.0 terms and additional component notices. Future adoption needs the selected distribution, its notices, dependency identities and an operational maintainer.~\citep{exec24_ray_sources}

The prior mutable README remains historical discovery evidence. The current review uses authenticated source blobs at the client cutoff. Where an acquired later-commit file has the same Git blob as the cutoff tree, that equality is recorded; five changed documentation files were acquired from the corrected cutoff itself. This prevents a convenient newer explanation from being silently substituted for the selected artifact. The profile focuses on Core tasks, actors, objects, resource declarations and runtime environments. The separate Train, Tune, Serve, Data and RLlib families are not admitted through Core's name. A commercial hosted service likewise requires its own entitlement, security, price and support evidence.~\citep{exec24_ray_sources,exec24_ray_contract}

\subsection{Retries do not establish exactly-once scientific accounting}
The task documentation distinguishes worker or machine failure from application exceptions. It describes three default retries for system failures, with explicit configuration and environment overrides; application exceptions are not retried by default. The inspected Python submission code fills defaults, reads the retry environment value and preserves explicitly supplied options. It serializes the function for export and can re-export it for a new driver context. A recovered task must therefore bind its input, code, closure, random seed and environment identities, not merely its display name.~\citep{exec24_ray_sources}

The object documentation describes reconstruction by rerunning the producer and recursively reconstructing arguments. It requires a live owner, suitable lineage and deterministic, idempotent tasks; objects created by the selected put interface do not have the same reconstruction route. The cutoff documentation's broad phrase about an environment setting disabling reconstruction must be read alongside the submission code's explicit-option precedence. This review makes only that narrower source-supported statement and leaves native reconstruction behavior under each configuration to conformance testing. A successful first return does not establish that the producer will never execute again.~\citep{exec24_ray_sources}

Let a scientific request require four logical contributions $z_A=2$, $z_B=4$, $z_C=-1$ and $z_D=3$ USD, so its mean is 2 USD. If a retry adds another contribution from $B$, averaging five attempts gives 2.40 USD, while keeping the original denominator of four gives 3 USD. Neither is the requested mean. An authored reducer keyed by immutable contribution identity yields 2 USD for all 120 permutations of the five events. If a duplicate identity carries a different value, it refuses the conflicting result instead of selecting whichever arrived last. This finite witness is not a Ray fault experiment; it specifies behavior a future adapter must demonstrate.~\citep{exec24_ray_contract}

The contribution identity should include the study, dataset snapshot, cluster or fold, estimator configuration, policy target and code/environment digest. Attempt identity is separate. An external payment, audit label purchase or model-service call needs its own idempotency protocol: recording a result after a side effect does not prove the effect occurred once. If an acknowledgement is lost, the client must represent an unresolved outcome until reconciliation establishes it. An MCP retry and a runtime retry must not multiply the same logical request invisibly.

\subsection{Actor recovery, object ownership and durable state}
The selected actor documentation gives no automatic restart by default. Configured restart reruns the constructor; it does not restore application state automatically. With retries disabled, an error can still arrive after a method executed. Enabling actor-task retries permits repeated execution, so an accumulator, optimizer or budget ledger needs an application checkpoint and recovery protocol. Async and threaded ordering also has distinct caveats. A checkpoint on an unavailable failed node is not enough when the actor restarts elsewhere.~\citep{exec24_ray_sources}

Object ownership is separate from the worker that computes a value. The selected documentation does not promise recovery after object-owner failure. Local object spilling extends capacity but is not a durable study ledger. At the cluster layer, the GCS stores metadata in memory by default; the reviewed text distinguishes external Redis support in a specified KubeRay/Serve setting from an alpha embedded RocksDB alternative and identifies functions unavailable during recovery. These mechanisms do not restore a bank's scientific history, actor application state or external side effects by implication. Full qualification must test driver, owner, worker, head-node and durable-store failures separately.~\citep{exec24_ray_sources}

A practical design stores immutable manifests and completed scientific contributions in an independently recoverable store, with atomic acceptance of a contribution identity. Workers compute candidates; a coordinator reconciles them against the required manifest. Recovery restores the manifest, verifies accepted content and reruns only unresolved logical work. A failed or timed-out request remains in the denominator even when a later retry succeeds. This design is proposed here; no storage service or Ray cluster has been deployed.~\citep{exec24_ray_contract}

\subsection{Parallel partitions are not statistical sampling units}
A distribution scheme must preserve both the estimator and its uncertainty calculation. For four row scores $(1,1,-1,-1)$, suppose the first two rows belong to one household and the other two to another. The uncorrected cluster-meat quantity normalized by the square of four is
\[
 \frac{(1+1)^2+(-1-1)^2}{4^2}=\frac12.
\]
Replacing household sums by row squares gives $1/4$. Splitting a household across workers and summing local squares loses the cross terms. Neither number by itself supplies a calibrated confidence interval; the example isolates a deterministic aggregation error before any sampling-law argument. A correct implementation combines fragments by the declared sampling unit before applying the covariance construction.~\citep{exec24_ray_contract}

Training and evaluation folds similarly attach to stable unit identities and the information clock, not worker number or arrival order. The authored fold map hashes a fixed study key and household identity; repeated rows of the household retain their fold under a row permutation. Assigning alternating rows instead splits each household in the example. Time boundaries and common calendar shocks need their own restrictions; hashing is not a cure for dependence or future-information leakage. Retry seeds should derive from the logical task identity, while deliberately repeated Monte Carlo replications receive distinct registered replication identities.

Even after identity preservation, floating-point reductions can change with order. In the authored binary-floating-point example, $(10^{16}-10^{16})+1$ is 1, while $10^{16}+(-10^{16}+1)$ is 0. A stable summation reference gives 1. The adapter should specify deterministic reduction or a justified numerical tolerance relative to the scientific error budget. Matching row counts and returning finite numbers is insufficient near a decision threshold. Exact equality can be a useful small-oracle test without being an appropriate universal requirement for every accelerator kernel.~\citep{exec24_ray_contract}

\subsection{Runtime environments and resource requests are operational contracts}
The pinned documentation distinguishes a job-submission environment, which applies before its entrypoint, from an environment supplied in initialization, which applies to child tasks and actors but not the driver. It describes inheritance, environment-variable merging, replacement of other fields, conflicts and an explicit override mode. It also warns that communicating tasks with incompatible package environments can fail during deserialization. A path-based local URI identifies whatever bytes are present when each worker starts, not an immutable code snapshot. These rules require a complete driver/worker environment manifest and actual artifact hashes.~\citep{exec24_ray_sources}

Runtime installation and cache behavior affect both availability and cost. Setup timeouts, cache misses, cold workers and repeated environment creation belong in the request budget. A runtime environment is not a security isolation guarantee. The source security model expects trusted code and external network and workload isolation; developer services can provide broad code-execution access. The pinned token-authentication documentation supplies an additional access-control mechanism, with rollout and configuration details that must be checked against the chosen distribution. Tokens do not replace transport encryption, tenant isolation or a narrow application authorization policy. No token, credential or live endpoint was accessed in this review.~\citep{exec24_ray_sources}

Logical CPU and memory resources govern admission, not physical enforcement. A task requesting one CPU can still launch multiple threads; GPU visibility and operating-system limits are separate mechanisms. Zero-copy NumPy access applies to read-only arrays on the same node, while mutation requires a copy. The selected documentation also warns about an optional shared-memory tensor path. Neither a resource request nor a zero-copy label proves a bank memory or isolation service level. The prototype must measure peak physical memory, serialization, network transfer, spill storage and process/thread behavior under its actual workload.~\citep{exec24_ray_sources}

\subsection{A bounded adoption decision and its economics}
The retained disposition is reference, with native integration deferred. A single-process or local worker pool is the baseline. Original Ray Core workers become a candidate only if measured workload partitioning and data movement justify them; building a new distributed scheduler or porting Ray wholesale is not justified by this proposal. The scientific owner specifies fold, target and covariance invariants; the adapter owner handles logical identities, retries and numerical reductions; the platform owner owns isolation and recovery; and the evaluation owner measures complete costs and failure denominators.~\citep{exec24_ray_contract}

An authored scaling calculation illustrates the gate. If one fifth of a workload is serial, four workers give an ideal speedup of $1/(1/5+(4/5)/4)=2.5$. Adding overhead equal to three tenths of the original runtime reduces this to $10/7$, about 1.43. These are stipulated proportions, not a Ray benchmark. Likewise, stipulated costs of 180 USD compute, 120 USD engineering and 60 USD storage/operations total 360 USD. Six attempts on four logical requests, of which three yield accepted artifacts, cost 60 USD per attempt, 90 USD per logical request and 120 USD per accepted artifact. A cheaper successful attempt is not evidence of lower full workflow cost.~\citep{exec24_ray_contract}

The proposed CLI/MCP boundary accepts named, approved computations over immutable artifact references, not arbitrary user-supplied code or unrestricted cluster credentials. It returns the logical request identity, attempt ledger, accepted contributions, refusals, unresolved effects, numerical checks and measured cost components. Public/generated fixtures should freeze household/time partitions, known outputs, duplicate/conflicting-result events, crash positions, cold environments and timeout budgets before a native study. The acceptance question is whether the selected implementation preserves scientific meaning and useful complete cost through these events. Until that evidence exists, distributed runtime support is a conditional engineering route, not a qualified causal capability or a bank deployment.~\citep{exec24_ray_contract}

The choice among these platforms therefore turns on a specified workload and accepted scientific output. Existing commercial services may already provide dispatch and recovery; open clients offer narrower integration boundaries; local workers provide the simplest comparison where distribution has no demonstrated benefit. None removes the need to preserve the target, information clock, missing-result denominator and output identity. The next treatment follows those obligations inside the numerical runtime, where covariance order and uncertainty metadata can change an answer even after the job has completed correctly.

